\newpage
\section*{Задание 2: Реверс-инжиниринг промптов}
\subsection{Постановка задачи варианта}
\noindent
Референсный объект: Сгоревший серверный блок питания \\
Окружение / Контекст: Стол в мастерской \\
Стиль референса: Реалистичное фото последствий аварии \\
Сложность интерпретации: Передача текстуры копоти и оплавленного пластика \\
Критический визуальный маркер: Обугленные провода и следы гари на металлическом корпусе \\
\begin{figure}[H]
	\centering
	\includegraphics[width=17cm]{z2/0.png}
	\caption{Референс}
\end{figure}
\newpage
\subsection{Гипотеза и декомпозиция сцены}
Для получения эталонного референса был сгенерирован следующий промпт (без фиксации Seed, но с акцентом на ключевые детали):
\textit{<<Реалистичное фото последствий аварии: сгоревший серверный блок питания лежит на столе в мастерской. Металлический корпус покрыт копотью и следами гари, видны обугленные провода, оплавленный пластик. Высокая детализация, крупный план, рассеянный естественный свет, резкие тени.>>}
Полученное изображение использовалось как неизвестный референс для реверс-инжиниринга.
\subsection{Журнал генерации и отладки}
\textbf{Итерация 1 (первичное описание от VLM)} \\
Запрос к мультимодальной модели: \textit{<<Ты – профессиональный промпт-инженер. Проанализируй прикреплённое изображение и составь подробное текстовое описание, пригодное для повторной генерации в нейросети Text-to-Image. Разбей описание на блоки: [Субъект], [Окружение], [Стиль], [Освещение], [Цветовая палитра]. Не используй художественных отступлений, только фактические визуальные детали.>>} \\
\newline
\indent \textbf{Ответ VLM (полученный промпт v1):} \\
\textit{<<[Субъект]
Компьютерный блок питания (PSU) модели Corsair CX600M со следами катастрофического термического повреждения. Корпус из матового серого металла сильно обгорел, покрыт черной сажей и окалиной. Пластиковая вставка в месте выхода кабелей оплавлена и деформирована. Из отверстия выходит пучок оплавленных черных проводов; изоляция кабелей слиплась, видны фрагменты оголенной меди и следы пузырящегося пластика. Коннекторы на концах кабелей частично разрушены огнем. На боковой грани видна узнаваемая этикетка с характеристиками и зеленым логотипом.
[Окружение]
Рабочий стол электронщика-ремонтника. На переднем плане — зеленый антистатический коврик для пайки с миллиметровой разметкой, засыпанный черным пеплом, обломками пластика и мелкими винтами. На заднем плане — коричневая перфорированная панель (пегборд) с аккуратно развешанными инструментами: отвертками с желто-черными рукоятками и плоскогубцами. Слева видна паяльная станция, катушка припоя и моток проводов. Справа — край материнской платы.
[Стиль]
Гиперреалистичная макрофотография с малой глубиной резкости (bokeh). Высокая детализация текстур: шершавый металл, глянцевый оплавленный пластик, ворсистая сажа. Композиция — ракурс три четверти, акцент на месте возгорания кабелей. Стиль промышленной съемки или документального фото инцидента.
[Освещение]
Контрастное, направленное освещение от настольной лампы на шарнире, расположенной сверху. Лампа создает яркое световое пятно на верхней крышке блока питания и глубокие тени в нижней части кадра. Теплый свет лампы (около 3000K) контрастирует с холодным дневным светом мастерской на заднем плане.
[Цветовая палитра]
Доминирующие цвета: угольно-черный, темно-серый (металл) и тускло-зеленый (коврик). Акценты: ярко-желтый и красный на рукоятках инструментов, белый и зеленый текст на наклейке блока питания. Цветовая гамма приглушенная, за исключением зоны под прямой засветкой лампы.>>}
\newpage
\textbf{Генерация по промпту v1:}
\begin{figure}[H]
	\centering
	\includegraphics[width=15cm]{z2/1a.png}
	\caption{Итерация 1, модель (a)}
\end{figure}
% \newpage
\begin{figure}[H]
	\centering
	\includegraphics[width=17cm]{z2/1b.png}
	\caption{Итерация 1, модель (b)}
\end{figure}
\newpage
\textbf{Анализ ошибок:}
Модель передала общую атмосферу и фактуру копоти, но провода выглядят хаотично. \\
\indent \textbf{Корректировка:}
Усилить описание проводов, добавив конкретные детали: <<провода выходят из одного круглого отверстия в углу>>, <<на блоке питания нет других отверстий>> \\
\newline
\indent \textbf{Итерация 2} \\
Промпт v2 (ручная доработка):
\textit{<<[Субъект]
Компьютерный блок питания (PSU) модели Corsair CX600M со следами катастрофического термического повреждения. Корпус из матового серого металла сильно обгорел, покрыт черной сажей и окалиной. Пучок оплавленных черных проводов выходит из одного круглого отверстия в углу с пластиковой вставкой. пластиковая вставка в месте выхода кабелей оплавлена и деформирована. На блоке питания нет других отверстий или контактов для подключения. Изоляция кабелей слиплась, видны фрагменты оголенной меди и следы пузырящегося пластика. Коннекторы на концах кабелей частично разрушены огнем. На боковой грани видна узнаваемая этикетка с характеристиками и зеленым логотипом.
[Окружение]
Рабочий стол электронщика-ремонтника. На переднем плане — зеленый антистатический коврик для пайки с миллиметровой разметкой, засыпанный черным пеплом, обломками пластика и мелкими винтами. На заднем плане — коричневая перфорированная панель (пегборд) с аккуратно развешанными инструментами: отвертками с желто-черными рукоятками и плоскогубцами. Слева видна паяльная станция, катушка припоя и моток проводов. Справа — край материнской платы.
[Стиль]
Гиперреалистичная макрофотография с малой глубиной резкости (bokeh). Высокая детализация текстур: шершавый металл, глянцевый оплавленный пластик, ворсистая сажа. Композиция — ракурс три четверти, акцент на месте возгорания кабелей. Стиль промышленной съемки или документального фото инцидента.
[Освещение]
Контрастное, направленное освещение от настольной лампы на шарнире, расположенной сверху. Лампа создает яркое световое пятно на верхней крышке блока питания и глубокие тени в нижней части кадра. Теплый свет лампы (около 3000K) контрастирует с холодным дневным светом мастерской на заднем плане.
[Цветовая палитра]
Доминирующие цвета: угольно-черный, темно-серый (металл) и тускло-зеленый (коврик). Акценты: ярко-желтый и красный на рукоятках инструментов, белый и зеленый текст на наклейке блока питания. Цветовая гамма приглушенная, за исключением зоны под прямой засветкой лампы.>>} \\
\newline
\textbf{Генерация по промпту v2:}
\begin{figure}[H]
	\centering
	\includegraphics[width=17cm]{z2/2a.png}
	\caption{Итерация 2, модель (a)}
\end{figure}
\newpage
\begin{figure}[H]
	\centering
	\includegraphics[width=17cm]{z2/2b.png}
	\caption{Итерация 2, модель (b)}
\end{figure}
% \newpage
\textbf{Анализ ошибок:}
Уже ближе к оригиналу, но блок питания должен быть светло-серым, а освещение белое от потолочной лампы и жёлтоватое от настольной лампы. \\
\indent \textbf{Корректировка:}
Добавить в промпт корректировку по цвету блока питания и освещению в комнате. \\
\newline
\indent \textbf{Итерация 3 (финальная)} \\
Промпт v3 (ручная доработка):
\textit{<<[Субъект]
Компьютерный блок питания (PSU) модели Corsair CX600M со следами катастрофического термического повреждения. Корпус из матового строго светло-серого металла сильно обгорел, покрыт черной сажей и окалиной. Пучок оплавленных черных проводов выходит из одного круглого отверстия в углу с пластиковой вставкой. пластиковая вставка в месте выхода кабелей оплавлена и деформирована. На блоке питания нет других отверстий или контактов для подключения. Изоляция кабелей слиплась, видны фрагменты оголенной меди и следы пузырящегося пластика. Коннекторы на концах кабелей частично разрушены огнем. На боковой грани видна узнаваемая этикетка с характеристиками и зеленым логотипом.
[Окружение]
Рабочий стол электронщика-ремонтника. На переднем плане — зеленый антистатический коврик для пайки с миллиметровой разметкой, засыпанный черным пеплом, обломками пластика и мелкими винтами. На заднем плане — коричневая перфорированная панель (пегборд) с аккуратно развешанными инструментами: отвертками с желто-черными рукоятками и плоскогубцами. Слева видна паяльная станция, катушка припоя и моток проводов. Справа — край материнской платы.
[Стиль]
Гиперреалистичная макрофотография с малой глубиной резкости (bokeh). Высокая детализация текстур: шершавый металл, глянцевый оплавленный пластик, ворсистая сажа. Композиция — ракурс три четверти, акцент на месте возгорания кабелей. Стиль промышленной съемки или документального фото инцидента.
[Освещение]
Строго: Контрастное, направленное освещение от настольной лампы на шарнире, расположенной сверху. Лампа создает яркое световое пятно на верхней крышке блока питания и глубокие тени в нижней части кадра. Теплый свет лампы (около 3000K) контрастирует с строго холодным белым светом мастерской от потолочной лампы на заднем плане.
[Цветовая палитра]
Доминирующие цвета: угольно-черный, темно-серый (металл) и тускло-зеленый (коврик). Акценты: ярко-желтый и красный на рукоятках инструментов, белый и зеленый текст на наклейке блока питания. Цветовая гамма приглушенная, за исключением зоны под прямой засветкой лампы.>>} \\
\newpage
\textbf{Генерация по промпту v3:}
\begin{figure}[H]
	\centering
	\includegraphics[width=17cm]{z2/3a.png}
	\caption{Итерация 3, модель (a)}
\end{figure}
\newpage
\begin{figure}[H]
	\centering
	\includegraphics[width=17cm]{z2/3b.png}
	\caption{Итерация 3, модель (b)}
\end{figure}
% \newpage
\textbf{Анализ:} \\
Достигнуто максимальное сходство с оригиналом. Ключевые маркеры (обугленные провода, оплавленный пластик, следы гари на металле) воспроизведены корректно. Семантический разрыв преодолён за счёт многократной детализации и языкового управления.
\newpage
\subsection{Визуализация результатов}
\begin{figure}[H]
	\centering
	\includegraphics[width=17cm]{z2/0.png}
	\caption{Референс}
\end{figure}
\begin{figure}[H]
	\centering
	\includegraphics[width=17cm]{z2/3b.png}
	\caption{Итерация 3, модель (b)}
\end{figure}
\newpage
Прямое описание от VLM (v1) дало достаточное приближённый результат, но некоторые детали колебались.
Последующие итерации с ручным уточнением позволили добиться высокой степени сходства,
особенно важным оказалось добавление анатомически верного расположения проводов и уточнение материалов.
\subsection{Сравнение моделей}
\begin{table}[h!]
	\centering
	\begin{tabular}{|l|p{5cm}|p{5cm}|}
		\hline
		\textbf{Критерий} & \textbf{модель (a)} & \textbf{модель (b)} \\
		\hline
		Следование промпту & 4 & 5 \\
		\hline
		Качество деталей & 3 & 5 \\
		\hline
		\textbf{Итог} & 3.5 & 5\\
		\hline
	\end{tabular}
\end{table}
Модель (b) справилась лучше, так как в меру следует промпту даже без параметров, в свою очередь модель (a) хуже справилась с генерацией так как имеет меньшее количество шагов и хуже следует промпту.
\subsection{Финальный промпт}
Итоговый промпт: \textit{<<[Субъект]
Компьютерный блок питания (PSU) модели Corsair CX600M со следами катастрофического термического повреждения. Корпус из матового строго светло-серого металла сильно обгорел, покрыт черной сажей и окалиной. Пучок оплавленных черных проводов выходит из одного круглого отверстия в углу с пластиковой вставкой. пластиковая вставка в месте выхода кабелей оплавлена и деформирована. На блоке питания нет других отверстий или контактов для подключения. Изоляция кабелей слиплась, видны фрагменты оголенной меди и следы пузырящегося пластика. Коннекторы на концах кабелей частично разрушены огнем. На боковой грани видна узнаваемая этикетка с характеристиками и зеленым логотипом.
[Окружение]
Рабочий стол электронщика-ремонтника. На переднем плане — зеленый антистатический коврик для пайки с миллиметровой разметкой, засыпанный черным пеплом, обломками пластика и мелкими винтами. На заднем плане — коричневая перфорированная панель (пегборд) с аккуратно развешанными инструментами: отвертками с желто-черными рукоятками и плоскогубцами. Слева видна паяльная станция, катушка припоя и моток проводов. Справа — край материнской платы.
[Стиль]
Гиперреалистичная макрофотография с малой глубиной резкости (bokeh). Высокая детализация текстур: шершавый металл, глянцевый оплавленный пластик, ворсистая сажа. Композиция — ракурс три четверти, акцент на месте возгорания кабелей. Стиль промышленной съемки или документального фото инцидента.
[Освещение]
Строго: Контрастное, направленное освещение от настольной лампы на шарнире, расположенной сверху. Лампа создает яркое световое пятно на верхней крышке блока питания и глубокие тени в нижней части кадра. Теплый свет лампы (около 3000K) контрастирует с строго холодным белым светом мастерской от потолочной лампы на заднем плане.
[Цветовая палитра]
Доминирующие цвета: угольно-черный, темно-серый (металл) и тускло-зеленый (коврик). Акценты: ярко-желтый и красный на рукоятках инструментов, белый и зеленый текст на наклейке блока питания. Цветовая гамма приглушенная, за исключением зоны под прямой засветкой лампы.>>}
\subsection{Вывод}
Реверс-инжиниринг промпта потребовал итеративного подхода, так как прямое описание от VLM оказалось недостаточно детальным для воспроизведения сложных визуальных признаков. Ключевым приёмом стало использование языкового управления вместо числовых параметров: добавление модификаторов «строгое соответствие», «критически важно», конкретизация местоположения объектов («из отверстия в левом углу, нет других отверстий») позволило компенсировать отсутствие доступа к Seed и CFG Scale. В итоге удалось добиться высокой степени сходства, применимой для задач инженерной визуализации.